\documentclass[12pt]{article}
\usepackage[utf8]{inputenc}
\usepackage[T1]{fontenc}
\usepackage{hyperref}
\title{Wanap Language Resources for K-12 Curriculum Development}
\author{Compiled by Claude Sonnet 5.5 (claude-sonnet-5-5) for Andr\'as Kornai}
\date{}
\begin{document}
\maketitle

\section{Introduction}
Wanap (ISO 639-3 wnp; also called Kayik) is a Torricelli language of the Palei branch, spoken in the mountainous eastern part of Sandaun Province, Papua New Guinea. Reported speaker numbers are about 1,100 (2003) on Wikipedia and 1,070 in the ASJP wordlist. Joshua Project gives a people-group population of about 2,100, which is not the same as a speaker count. Little educational or documentary material turned up in search, so this report is short.

\section{Primary Educational Resources}
No primers, textbooks or school programmes for Wanap were found in the search results. Joshua Project's language page lists no written form and says Bible translation is needed.

\section{Linguistic and Reference Materials}
\begin{itemize}
\item Wikipedia, Wanap language (overview and classification): \url{https://en.wikipedia.org/wiki/Wanap_language}
\item Wikipedia, Palei languages (the branch Wanap belongs to): \url{https://en.wikipedia.org/wiki/Palei_languages}
\item ASJP wordlist, Wanap (a short comparative word list only): \url{https://asjp.clld.org/languages/WANAP_2}
\end{itemize}
I found no grammar or dictionary in the search results.

\section{Religious and Community Resources}
\begin{itemize}
\item Joshua Project language entry for Wanap (reports a need for Bible translation and no written form): \url{https://joshuaproject.net/languages/wnp}
\item Joshua Project people-group page, Wanap in Papua New Guinea: \url{https://joshuaproject.net/index.php/people_groups/15850}
\end{itemize}
I found no scripture translations.

\section{Digital and Audio Resources}
No audio, video or corpus resources for Wanap were found in the search results.

\section{Curriculum Development Context}
Wanap is spoken by all adults but not by all young people, and it is not known to be taught in schools. Curriculum work would have to start with new documentation: community-led orthography development, recordings, and a basic word list. The sources found so far are only overview-level, so primary fieldwork with the community is needed.
\end{document}